%%%PAGE 8 rot=0 legibility=good summary=重核裂变与轻核聚变；半衰期（统计规律）；核力与结合能、比结合能曲线；氘氚聚变的能量与动量问题；光强-λ图与光电子；波粒二象性、电子波不能用牛顿定律、不确定关系
%%%BLOCK id=008-01 ch=17 sec=重核裂变 kind=knowledge alt=none cont=none y=0.04-0.11
\textbf{重核裂变}（链式反应）\quad 1中子$\to$多于2个中子\qquad 条件：体积足够大到临界体积

\[
\nuc{1}{0}{n}+\nuc{238}{92}{U}\to\mathrm{Sr}+\mathrm{Ba}+3\,\nuc{1}{0}{n}\qquad\text{原子弹、核电站}
\]
% CHECK: $\mathrm{U}$ 的质量数手写形似238（常见裂变用235）；第一个产物符号形似“Sr”（亦可能是Kr），均照原样转录
%%%END
%%%BLOCK id=008-02 ch=17 sec=轻核聚变 kind=knowledge alt=none cont=none y=0.11-0.19
\textbf{轻核聚变}\quad 两小生一大，中子是附加\unclear{品}\qquad 条件：几百万度高温，热核反应

\[
\nuc{2}{1}{H}+\nuc{3}{1}{H}\to\nuc{4}{2}{He}+\nuc{1}{0}{n}\qquad\text{（用原子弹引爆）}\quad\text{也叫轻核骤变}
\]
% CHECK: “骤变”疑为“聚变”（原稿字形为“骤”）
%%%END
%%%BLOCK id=008-03 ch=17 sec=半衰期 kind=knowledge alt=none cont=none y=0.22-0.36
\textbf{半衰期}\quad 未衰变物质质量\ $m=M_0\left(\dfrac{1}{2}\right)^n$\qquad $n$个半衰期

\hspace{4.5em}已衰变物质质量\ $M_0-m$\qquad\qquad 说克不说个，是大量

\unclear{衡}量衰变快慢的物理量\qquad 16个\unclear{氡}经1个半衰期，还剩8个\ $\times$

只由元素决定\qquad 16g\unclear{氡}经1个半衰期，还剩$\underset{\approx16\unit{g}}{8\unit{g}}$\ $\times$\quad 还剩$8\unit{g}$未衰变
% 原稿“还剩8g”下方有小字“≈16g”；“氡”字（疑）原稿写成圈起来的“80”状字形；左侧两行“（衡）量衰变快慢的物理量”“只由元素决定”位于下方两行的左边
%%%END
%%%BLOCK id=008-04 ch=17 sec=核力与结合能 kind=knowledge alt=none cont=none y=0.38-0.62
\textbf{核力}\quad 核子与核子\ $\downarrow$\ 强力（\unclear{交换介子}）\ $\downarrow$ % 原稿竖排：“核子与核子”下接箭头指向“强力（…）”，其下再有一箭头指向“结合能”；括号内字被箭头竖线穿过，难辨

\begin{itemize}
\item 核子\fbox{分开}强力做负功，\fbox{吸收}能量，$m$增大
\item 核子\fbox{结合}强力做正功，\fbox{放出}能量，$m$减少
\end{itemize}
% 以上两行原稿共用一个大括号（括号左侧即“强力（…）”）；加框的词在原稿中是手绘圆圈圈出

%FIG p008_a 0.66 0.425 0.88 0.545
\notefig[0.3]{figures/p008_a.png}

结合能$=$比结合能$\times$核子数

（键能）\quad $\Delta E=\Delta m\cdot c^2$

核子数越多，结合能越大\quad Fe % “Fe”小字写在这一行末尾右侧上方（对应上图曲线峰值处的Fe）

比结合能（结合能/核子数量）：比结合能越大，原子核越稳定
%%%END
%%%BLOCK id=008-05 ch=17 sec=核能·质能方程 kind=knowledge alt=none cont=none y=0.69-0.90
\[
\nuc{2}{1}{H}+\nuc{3}{1}{H}\to\nuc{4}{2}{He}+\nuc{1}{0}{n}+17.6\unit{MeV}
\]

%FIG p008_b 0.05 0.738 0.31 0.79
\notefig[0.27]{figures/p008_b.png}

$1.09\unit{MeV}$\quad $2.78\unit{MeV}$ % 原稿写在上式两个反应物的下方，并带两条弧形箭头（见图 p008_b；图内已含此二数）

\[
E_{\text{前}}+\underset{17.6\unit{MeV}}{E_{\text{放}}}=E_{\text{后}}\qquad Q=\Delta mc^2
\]

若\unclear{$\alpha$}以$Q$轰击$\nuc{14}{7}{N}$则不会反应，还需提供动能 % “若”后第一个字（形似“X”）疑为α

生成物总结合能$>$反应物总结合能\ \textsuperscript{\unclear{光子有动质量}}

新核和$\alpha$粒子动量不等大反向\quad 光子也有动量\quad 释放能量给了$\gamma$光子和新核的动能

不能说质量就是能量，不能表示质量和能量间转化关系\quad 可说正比
%%%END
%%%BLOCK id=008-06 ch=17 sec=光电效应 kind=knowledge alt=none cont=none y=0.89-1.00
%FIG p008_c 0.05 0.893 0.2265 0.985
\notefig[0.25]{figures/p008_c.png}

只看最小$\lambda$ % 图 p008_c 内亦含此字

$\lambda<\lambda_0$就可以产生光电子
%%%END
%%%BLOCK id=008-07 ch=17 sec=波粒二象性·物质波 kind=knowledge alt=none cont=none y=0.89-1.00
波粒二象性：$P=\dfrac{h}{\lambda}$\qquad $\lambda=\dfrac{h}{p}=\dfrac{h}{\sqrt{2mUe}}$

\begin{itemize}
\item 波动\quad 多\quad \unclear{概率}
\item 运动\quad 少\quad \unclear{相关}
\end{itemize}
% 原稿：“波粒二象性”下方两条箭头分别指向“波动”与“运动”；“波动”下有小字“多”及两字（疑“概率”）；“运动”下有小字“少”及“相关”

$\Delta x\,\Delta p\ge\dfrac{h}{4\pi}$\quad 概率波

电子波不能用牛顿第二定律，（无轨迹，无坐标）\ \red{（微观高\unclear{速}…} % 红字右侧被照片边缘截断

非经典力学问题\ \red{无法用牛顿运动定律}
%%%END
%%%PAGEEND 8
